\label{chap.4}
\section{Design Principles}
It is easy to score well on a benchmark by flagging anything that looks a bit suspicious. To rule that out, every attack in MAPIS-Bench comes with harmless \emph{twins} that look almost the same:
\begin{itemize}
\item \textbf{authorized}: the user named the destination or rule that the attack plants;
\item \textbf{resisted}: the planted fragments are present in the data, but the agent acts correctly;
\item \textbf{designated}: the user named the file to use.
\end{itemize}
Training and test data are split by attack family and wording, and the test families use action tools that never appear in training.

\section{The Six Sets}
\begin{table}[H]
\centering\small
\begin{tabular}{p{3.1cm}p{9.6cm}}
\toprule
Set & Content\\
\midrule
MAPIS-Bench v1 & 1,693 sessions: AgentDojo (1,000 attack and 132 benign) and InjecAgent (561 data-stealing cases), audited against the original sources. Splits 1,169 / 242 / 282.\\
BIPIA & 1,800 minimal-pair sessions. The BIPIA test part is used only for testing and was never trained on.\\
MAPIS-Decomposed & 288 sessions from 12 template families: exfiltration through memory where no single message looks harmful (evaluation only).\\
MAPIS-MultiHop & About 2,100 sessions. Attacks are fragments (a changed detail plus an authority claim) inside real carrier text across untrusted hops; the agent's final action uses them. Four skeletons realise the five mechanisms of the abstract; depth padding of 0, 2 or 4 benign hops.\\
MAPIS-RealHarm & 2,330 sessions from InjecAgent direct harm (510 cases) and ASB (10 agents, 400 attacker tools), with held-out attacker tools and agents.\\
MAPIS-Adaptive & 1,692 sessions of an attacker who knows MAPIS: obfuscated items, items split over two documents, invisible Unicode, paraphrased claims, delay beyond the four-hop window, and combinations.\\
\bottomrule
\end{tabular}
\caption{The benchmark sets.}
\end{table}

\section{Attack Classes and Mechanisms}
We report results in two ways: by the five attack classes from the first review (data exfiltration, financial manipulation, code/tool manipulation, instruction override, physical safety harm) and by the five mechanisms from the abstract (document poisoning, memory injection, tool-output manipulation, cross-agent instruction injection, goal hijacking). The four MultiHop skeletons are tool fragments (covering document poisoning and tool-output manipulation), memory promotion (memory injection), cross-agent relay, and goal hijacking.

\section{Labelling Policy}
Labels are given per event. A fragment event on its own is labelled safe, because it looks the same in the attack, the resisted twin and the authorized twin, and the agent has not done anything yet. A relay, plan or memory-write message is labelled malicious only if it carries a planted item that can be acted on. The final action is labelled malicious. We arrived at this policy after the two failures described in Section~\ref{sec:history}.

\section{Final Test Set}
The test split has 2,304 sessions: 576 from MAPIS-Adaptive (192 attacks, 384 twins), 234 from MAPIS-MultiHop (65 attacks, 169 twins), 364 RealHarm sessions with held-out attacker tools and agents, 282 from v1, 800 from BIPIA and 48 decomposed sessions. Every system is given exactly the same events.

We also use two sets only for testing, reported in Section~\ref{sec:harder}. MAPIS-StateAttack has 168 sessions, including 56 attacks aimed at the session record itself: evidence buried under filler hops, claims split into pieces, items introduced by an agent, items the agent has to decode, and items passed through memory. The second set has 20 attacks and 20 twins written in plain English by the four of us, without using the generators.

\section{Limits of the Data}
MultiHop and Adaptive are generated from templates, even though the families, wordings and tools in the test split are held out and the carrier text is real. They are probably easier than real traffic, which is why we added the two test-only sets.
